\section{模型假设与符号说明}

\subsection{基本假设}

\begin{enumerate}
  \item 附件 CSV 中的收盘价是本文拟合、校准和误差评价的唯一真实价格序列。
  \item 供应中断、SPR 释放上限、绕道运输能力等变量遵循赛题给定范围，不在范围外任意调整。
  \item 短期需求弹性按题面给定低弹性口径处理；中长期需求弹性允许逐步过渡，以反映高价下需求调整更充分。
  \item 模型不使用新闻爬虫数据，也不把外部新闻情绪作为观测输入；恐慌强度和衰减率是行为机制参数。
  \item 60--180 天预测为条件情景外推，不代表真实未来观测，也不编造任何未来真实价格。
\end{enumerate}

\subsection{符号说明}

\begin{table}[H]
\centering
\caption{主要符号说明}
\begin{tabularx}{\linewidth}{lXl}
\toprule
符号 & 含义 & 单位 \\
\midrule
$P_t$ & 第 $t$ 日实际收盘价 & 美元/桶 \\
$\hat P_t$ & 第 $t$ 日模型递推价格 & 美元/桶 \\
$P_0$ & 冲突窗口前收盘基准价 & 美元/桶 \\
$S_0,D_0$ & 战前基准供给和需求 & 万桶/日 \\
$I$ & 封锁造成的供应中断量 & 万桶/日 \\
$A_t$ & 其他产油国增产或替代供应增量 & 万桶/日 \\
$R_t$ & 第 $t$ 日 SPR 释放量 & 万桶/日 \\
$Q_t$ & 第 $t$ 日绕道运输恢复量 & 万桶/日 \\
$B_t$ & 商业库存当日缓冲量 & 万桶/日 \\
$G_t$ & 剩余供需缺口 & 万桶/日 \\
$\varepsilon_t$ & 需求价格弹性 & 无量纲 \\
$F_t$ & 恐慌溢价强度 & 无量纲 \\
$\beta$ & 地缘风险权重 & 无量纲 \\
$U$ & 不确定性与制度风险强度 & 无量纲 \\
$v,v_L$ & 短期和长期价格调整速度 & 无量纲 \\
$J_t$ & 长期库存耗尽或二次跳涨风险项 & 美元/桶 \\
\bottomrule
\end{tabularx}
\end{table}

\subsection{参数来源与识别边界}

为了避免把模型理解为单纯曲线拟合，本文将参数按来源分为六类。题面参数和附件统计量构成模型硬约束；官方外部数据只用于数量级审计；校准参数必须通过基准对比、消融实验和扰动检验；长期情景参数只用于 60--180 天条件外推，不参与短期窗口的误差拟合；长期机制强度参数则必须显式命名，并进入敏感性分析。

\begin{table}[H]
\centering
\caption{参数来源分类与使用边界}
\begin{tabularx}{\linewidth}{lX X}
\toprule
参数类别 & 代表变量 & 使用边界 \\
\midrule
题面给定参数 & 短期弹性 \(\varepsilon\)、供应中断范围、SPR 释放上限、恐慌持续时间 & 作为模型约束和传统供需基准依据，不在题面范围外任意调整 \\
附件统计量 & 基准价格 \(P_0\)、冲突窗口真实收盘价、收益率和波动率 & 作为拟合目标、误差评价和历史窗口比较依据 \\
官方外部审计量 & EIA 周度 SPR、商业库存、原油产量，JODI 多国月度石油数据，Caldara-Iacoviello GPR 指数 & 只用于审计长期约束、数量级和风险高位事实，不替代附件真实价格 \\
机制校准参数 & 价格传导系数 \(\phi\)、风险权重 \(\beta\)、缓冲确认强度、预期修复强度、调整速度 \(v\) & 在固定机制结构内校准，并接受消融、扰动和硬编码审计 \\
长期情景参数 & OPEC+ 响应、非 OPEC 供给恢复、长期弹性、制度风险强度 & 只用于条件中心路径和尾部风险外推，不称为真实未来拟合 \\
长期机制强度参数 & 冲击不确定性占比、制度风险占比、信心恢复后风险底座、过剩供给回归强度、封锁风险衰减速度 & 用于刻画长期机制形状，不作为官方物理常数，并在敏感性分析中单独扰动检验 \\
\bottomrule
\end{tabularx}
\end{table}

其中最容易被误读的是价格传导系数 \(\phi\)。本文没有用 \(\phi\) 替换题面短期弹性 \(\varepsilon=-0.05\)：低弹性仍用于传统供需反事实上界，解释为什么静态模型会给出 278--337 美元/桶的高价压力；\(\phi\) 只表示在 SPR、库存、绕道运输、需求收缩和预期修复共同存在时，剩余物理缺口进入布伦特期货目标价格的传导比例。二者分别对应“无缓冲静态供需压力”和“有缓冲期货定价传导”，经济含义不同。

更具体地说，若令 \(\phi=1\)，则模型等价于假设剩余物理缺口会按题面低弹性完全、即时地映射为期货目标价格；但现实原油期货市场同时定价库存、预期、投机头寸、政策缓冲可信度和未来需求调整。Kilian（2009）、Hamilton（2009）以及 Kilian 与 Murphy（2014）均指出，油价冲击不能只由当期物理供给缺口解释，预期需求、库存和投机交易会改变现货与期货价格的传导路径。因此本文把 \(\phi\) 明确归类为行为校准参数，并在反驳性检验中单独报告：当 \(\phi=1\) 时，短期 RMSE 升至 26.55 美元/桶，模拟峰值升至 166.84 美元/桶，说明完全机械传导不符合附件观测；而在 \(0.5\phi\)--\(2.0\phi\) 扰动内，RMSE 仍处于 3.50--3.64 美元/桶区间，说明价格平台并非唯一参数点偶然形成。

长期外推中的若干机制系数也需要单独说明。例如冲击不确定性占比、制度风险占比、过剩供给回归强度和封锁风险衰减速度并非 EIA、JODI 或 IEA 直接发布的物理量，而是把短期校准后的风险定价机制延伸到 60--180 天时使用的形状参数。本文不把它们写成确定事实，而是把它们显式列入敏感性分析：若这些系数在合理范围内扰动后结论仍保持中性回落、悲观尾部上行的结构，则说明长期判断主要来自机制组合，而不是某一个隐藏常数。
